% CHECKED DEDUCTION: complete self-read and archive independent audit passed.
% Include after greenberg_kernel_section.tex and integral_classification_section.tex.
\subsection{Integral odd generation under the one-variable Greenberg hypothesis}
\label{sec:greenberg-integral}

The integral classification extends from Vandiver to the weaker hypothesis
used in Theorem~\ref{thm:greenberg-kernel}.  The additional arithmetic
point is that irregularity already forces a nonzero minus class group;
Vandiver is unnecessary for that implication.

\begin{corollary}[K17-GREENBERG-INTEGRAL]
\label{cor:greenberg-integral}
Let $p$ be an odd prime, let
\[
 X=\varprojlim_r\operatorname{Cl}(\mathbb Q(\mu_{p^r}))\{p\}
\]
with norm transition maps, and assume that $X^+$ is finite.
Let $\mathfrak g_p$ be the actual integral tripod Lie algebra with
its arithmetic weight grading.
If $p$ is regular, then $\mathfrak g_p$ is the free
$\mathbb Z_p$-Lie algebra on one generator in each odd weight
$m\ge3$.  If $p$ is irregular, then
\[
 \bigl(\mathfrak g_p/[\mathfrak g_p,\mathfrak g_p]\bigr)_n
\]
is a nonzero finite $p$-group for infinitely many even weights $n$.
In particular, under the stated hypothesis, homogeneous odd integral
generators exist exactly at the regular odd primes, and at those primes
they can be chosen freely.
\end{corollary}

\begin{proof}
The regular-prime assertion is the one in
Theorem~\ref{thm:integral-classification}; its positive proof uses
Sharifi's Theorem~1.1 and Brown's rational Deligne--Ihara theorem
\cite[Theorem~1.1]{Sharifi2002}\cite[\S\S2.1--2.3]{Brown2019}.
For an irregular prime, put $F_\infty=\mathbb Q(\mu_{p^\infty})$,
\[
 H=\operatorname{Gal}(\Omega_p/F_\infty),\qquad
 M=H/\overline{[H,H]},\qquad
 A_0=\operatorname{Cl}(\mathbb Q(\mu_p))\{p\}.
\]
Theorem~\ref{thm:greenberg-kernel} identifies $H$ faithfully with
the actual cyclotomic-kernel image.  Its integral weight filtration
$F^nH$ is separated, and Proposition~\ref{prop:outer-sat}
identifies it with the actual outer filtration.

Here is the precise nonvanishing input.  Kummer's criterion says that
irregularity is equivalent to $p$ dividing the numerator of some
$B_k$ with $2\le k\le p-3$ even.  Ribet's Main Theorem~(1.1) gives
\[
 p\mid B_k
 \quad\Longleftrightarrow\quad
 (A_0/pA_0)^{(\chi^{1-k})}\ne0,
 \tag{GI.1}
\]
where $\chi$ is the mod-$p$ cyclotomic character
\cite[\S1, p.~151, Main Theorem~(1.1)]{Ribet1976}.
Since $1-k$ is odd, complex conjugation acts as $-1$ on this
eigenspace.  Thus $A_0^-\ne0$.  Conversely $A_0^-\ne0$ implies
that $p$ divides the full cyclotomic class number, hence that $p$
is irregular.  This is the exact irregularity/minus-class criterion
used here; it does not require $A_0^+=0$.

Let $A_\infty$ be the direct limit of the layer $p$-class groups
under extension of ideals.  The established Kummer comparisons and
minus-part injectivity give
\[
 A_0^-\hookrightarrow A_\infty^-,\qquad
 M^+\cong(A_\infty^-)^\vee(1),\qquad
 M^+\sim(X^-)^\iota(1),
 \tag{GI.2}
\]
where $\sim$ denotes a pseudo-isomorphism and $\vee$ is
Pontryagin duality
\cite[Propositions~3.4.1 and~3.4.7, Corollary~3.4.8]{SharifiIwasawaNotes}.
As in the proof of Theorem~\ref{thm:integral-classification},
Iwasawa's freeness input makes $H$ free pro-$p$, so $M$ is
$\mathbb Z_p$-torsion free; Ferrero--Washington's $\mu=0$
theorem and (GI.2) make $M^+$ finitely generated over $\mathbb Z_p$.
The injection and actual duality in (GI.2) make $M^+$ nonzero.
Consequently it is finite free of positive rank.  These steps use
no Vandiver assertion.

For completeness, the filtration argument is as follows.  Let
$M^n$ be the image of $F^nH$ in $M$.  These are closed submodules
and their intersection is zero: for a class $x$ in the intersection,
the nested nonempty compact fibres
$F^nH\cap\pi^{-1}(x)$ have a common point in
$\bigcap_nF^nH=\{1\}$, so $x=0$.
There are surjections
\[
 (\mathfrak g_p^{\mathrm{ab}})_n
 \twoheadrightarrow M^n/M^{n+1}.
 \tag{GI.3}
\]
Complex conjugation acts by $(-1)^n$ on these graded quotients.
Since $2$ is a unit in $\mathbb Z_p$, taking its plus eigenspace
is exact.  The filtration of $M^+$ therefore has zero quotients
in all odd weights.

Each integral graded Lie piece is finite free over $\mathbb Z_p$,
by its embedding in the outer Johnson module.  Its degreewise
abelianization is finitely generated, and Brown's rational theorem
makes that abelianization zero after inverting $p$ in every even
weight.  Thus each even $(\mathfrak g_p^{\mathrm{ab}})_n$, and
each even filtered quotient of $M^+$ in (GI.3), is a finite
$p$-group.  If only finitely many of the latter quotients were
nonzero, the separated filtration of $M^+$ would eventually
be zero and would express $M^+$ as an iterated extension of
finitely many finite groups.  This contradicts its positive
free rank.  Infinitely many even quotients, and hence infinitely
many even Lie indecomposable pieces in (GI.3), are nonzero.
An odd homogeneous generating collection would make all even
Lie indecomposables zero, completing the proof.
\end{proof}

The conclusion is conditional on the one-variable Greenberg hypothesis
when Vandiver is not known.  It gives neither the first exceptional
weight nor its order, and it makes no additional assertion about
primitivity of a prescribed normalization of the full-domain
Soul\'e characters.
